%% P19 \dash{} Metoda gradientu prostego: od SGD do Adama
%%
%% Zalążek. Poniższy opis jest kontraktem tego programu: co ma uzasadnić i co
%% ma dać czytelnikowi. Napisanie programu oznacza usunięcie bloku
%% \programstub{} i zastąpienie go programem. Jeśli po przerobieniu materiału
%% opis okaże się błędny, popraw go i odnotuj sprzeczność w CLAUDE.md w sekcji
%% *Resolved questions*.

\program{Metoda gradientu prostego: od SGD do Adama}
\label{prog:P19}

\programstub{%
\textit{[Opis do przetłumaczenia \dash{} patrz wydanie angielskie.]}

Argument: every optimiser in common use is the same update with a different
estimate of \enquote{how far and in which direction}, and each addition fixes a
specific named failure of the one before it. Payoff: the step-size bound
from P16; momentum as the exponential moving average met in F4, fixing the
zig-zag from P14; per-coordinate scaling fixing badly scaled features; why
Adam divides by the square root of the second moment \dash{} it is a per-
coordinate estimate of the gradient's scale, so the update becomes
approximately unit-sized regardless of how large the gradient is, which is
also precisely why Adam is insensitive to loss scaling and why the epsilon
must sit outside the square root; bias correction derived from the
initialisation at zero; weight decay is not L2 regularisation once you
divide by a running scale, which is the entire content of AdamW; warmup and
cosine schedules described as what they do to the effective step rather than
as ritual.

\medskip
\textbf{Planowana długość:} 65 ramek.
}
